% ============================================================
% Chapter 7: Results
%   Section: Model Comparison
%     Subsection: The Pipeline against Other Ways of Doing Each Phase
% Included from chapters/07-results/06-model-comparison/model-comparison.tex
% Source: backend/data/evaluation/experiments/frameworks/report.md,
% metrics.json and resources.json (2026-11-11). Every figure below is
% copied from report.md; rebuild it before changing one.
% ============================================================

\subsection{The Pipeline against Other Ways of Doing Each Phase}
\label{subsec:results-frameworks}

The comparison of Subsection~\ref{subsec:benchmark-frameworks} was run
on 11~November 2026. It is a first reading: the samples are small, and
what each table cannot show is said beside it.

\paragraph{Topics, on the project's own articles.}
Table~\ref{tab:frameworks-topics} compares five ways of naming an
article's topic among the pipeline's 23. The hosted models lead, but
their intervals overlap the embedding classifier's, and many of these
labels were written by a model and none has been reviewed by a person:
agreement with them is not yet accuracy. The two local language models
are slower than the classifier and no more accurate.

\begin{table}[ht]
\centering
\small
\begin{tabular}{@{}lcccc@{}}
\toprule
56 articles, 23 topics & Top-1 & Gold in top 3 & Top-1 group & Median s \\
\midrule
Embedding classifier (the pipeline's) & 0.66 [0.54, 0.79] & 0.88 & 0.82 & 4.7 \\
\texttt{llama3.2:3b} & 0.52 [0.38, 0.64] & 0.82 & 0.61 & 13.8 \\
\texttt{qwen3.5:4b} & 0.64 [0.52, 0.77] & 0.96 & 0.80 & 24.3 \\
Claude Haiku & 0.73 [0.61, 0.84] & 0.93 & 0.86 & 2.9 \\
Claude Sonnet & 0.77 [0.66, 0.88] & 0.96 & 0.89 & 2.6 \\
\bottomrule
\end{tabular}
\caption[Topic classification on the project's articles]{Topic
classification on the labelled articles. The language models read the
title and the first 3{,}000 characters, the embedding classifier the
whole body. 95\% bootstrap intervals; \texttt{llama3.2:3b} gave 7
unusable answers and \texttt{qwen3.5:4b} one.}
\label{tab:frameworks-topics}
\end{table}

\paragraph{Topics, on a standard set.} On SIB-200
(Table~\ref{tab:frameworks-sib}) training a classifier on the embeddings
the pipeline already computes adds about 15 points over using them
untrained, which is the pipeline's method, and equals the hosted model
at a twentieth of its time per sentence, deterministically. SIB-200 has
seven topics and single sentences where the pipeline has 23 topics and
whole articles, so the table shows what training buys the method, not
how the application's classifier scores; training one for the
pipeline's topics needs a few hundred labelled articles, and 56 exist.

\begin{table}[ht]
\centering
\small
\begin{tabular}{@{}lccc@{}}
\toprule
SIB-200 test, 204 sentences & English & Spanish & s / sentence \\
\midrule
\texttt{bge-m3}, untrained (the pipeline's method) & 0.75 [0.68, 0.80] & 0.74 [0.68, 0.80] & 0.10--0.14 \\
\texttt{bge-m3} + logistic regression & 0.89 [0.85, 0.93] & 0.91 [0.87, 0.95] & 0.10--0.14 \\
TF-IDF + logistic regression & 0.74 [0.67, 0.80] & 0.72 [0.66, 0.78] & $<0.001$ \\
Claude Haiku & 0.88 [0.83, 0.92] & 0.88 [0.83, 0.93] & 2.0--2.4 \\
\bottomrule
\end{tabular}
\caption[Topic classification on SIB-200]{Accuracy on SIB-200, the same
sentences in both languages, with 95\% bootstrap intervals. The trained
contenders learn from the language's 701 training sentences (76~s for
the embedding classifier, embedding included). One sentence was refused
by the hosted model's safety filter and is counted wrong.}
\label{tab:frameworks-sib}
\end{table}

\paragraph{Entities, on a standard set.} GLiNER as the pipeline runs it
is the weakest of the three serious contenders
(Table~\ref{tab:frameworks-ner}), with a precision of 0.50 in English
and an organisation score of 0.37, part of it the sports tables of
CoNLL-2003, where a city names a team. The fine-tuned model leads, but
it was trained on these corpora's training sets, which neither GLiNER
nor the hosted model was given: it is the supervised ceiling rather than
a like-for-like rival. GLiNER also extracts products and events, which
that model cannot. The public test sets may be in the hosted model's
training data.

\begin{table}[ht]
\centering
\small
\begin{tabular}{@{}llccccc@{}}
\toprule
 & & Precision & Recall & F1 & F1, any type & s / sentence \\
\midrule
GLiNER (the pipeline's) & en & 0.50 & 0.71 & 0.59 & 0.73 & 0.08 \\
 & es & 0.71 & 0.75 & 0.73 & 0.78 & 0.14 \\
XLM-RoBERTa, fine-tuned & en & 0.89 & 0.89 & 0.89 & 0.94 & 0.04 \\
 & es & 0.91 & 0.92 & 0.91 & 0.96 & 0.08 \\
spaCy, small & en & 0.64 & 0.54 & 0.58 & 0.71 & 0.003 \\
 & es & 0.60 & 0.61 & 0.60 & 0.80 & 0.004 \\
Claude Haiku & en & 0.83 & 0.86 & 0.84 & 0.92 & 2.3 \\
 & es & 0.89 & 0.87 & 0.88 & 0.92 & 2.5 \\
\bottomrule
\end{tabular}
\caption[Entity recognition on CoNLL]{Persons, organisations and
locations on 300 sampled test sentences each of CoNLL-2003 (en) and
CoNLL-2002 (es), micro-averaged. ``Any type'' counts a name found under
another type.}
\label{tab:frameworks-ner}
\end{table}

\paragraph{Verification, on the same evidence.} On ten AVeriTeC claims
(Table~\ref{tab:frameworks-verifier}) the hosted model was right on five
and the baseline on four, with intervals some 60 points wide: no
difference is shown. The finding is the last row. Asked with no evidence
at all, the hosted model was right on seven, more than with the
evidence: these fact-checks of 2020--22 are, in all likelihood, in what
it was trained on. No tool reached the web, so the rule that no evidence
may postdate the claim held for what it read; it cannot hold for what it
remembers. Any score of such a model on this benchmark is therefore an
upper bound, and seven of the ten labels being \emph{Refuted}, always
answering so scores seven too. The baseline, deterministic on fixed
evidence, repeated its verdicts of 5~November on all ten.

\begin{table}[ht]
\centering
\small
\begin{tabular}{@{}lcccc@{}}
\toprule
AVeriTeC dev, 10 claims & Right & Wrong and definite & Not verified & Median s \\
\midrule
\texttt{llama3.2:3b}, replayed & 4 & 4 & 2 & 70.8 \\
Claude Sonnet, replayed & 5 & 2 & 3 & 37.6 \\
Claude Sonnet, no evidence & 7 & 1 & 2 & 4.0 \\
Always \emph{Refuted} & 7 & 3 & 0 & --- \\
\bottomrule
\end{tabular}
\caption[The verifier on the same evidence]{Two verifiers on the
evidence one run had kept for the first ten claims of AVeriTeC's
development split, and the hosted model with none.}
\label{tab:frameworks-verifier}
\end{table}

\paragraph{Review.} On the planted-defect texts the hosted model scored
every defective text below its clean one (detection 100\%, against 95\%
for \texttt{gemma4:e4b} and \texttt{qwen3.5:9b}; $+20$ points over
\texttt{llama3.2:3b}, $[+12, +28]$), named every defect, and took 10.4~s
a text against 79.2~s for the reviewer in use. It is also the harshest
on clean text (a mean fact-consistency score of 80.5 where the local
reviewer gives 100) and moves the scores a defect was not aimed at more
(16.3 points against 7.3).

\paragraph{Resources.} Table~\ref{tab:frameworks-resources} gives what
each model needs. The models built for one task hold between 0.3 and
1.1~GB once loaded; the local language models between 3.7 and 7.9~GB,
five to ten times as much, and on topics were slower and no more
accurate. The fine-tuned entity model has nearly twice GLiNER's
parameters and peaked lower in memory. The trained topic classifier adds
some seven thousand weights to an embedding model the pipeline already
holds. The hosted model needs nothing local, and cannot be sized.

\begin{table}[ht]
\centering
\small
\begin{tabular}{@{}lrrrrr@{}}
\toprule
 & Parameters & On disk & RAM loaded & RAM peak & Cores \\
\midrule
\texttt{bge-m3} (embeddings) & 568 M & 2.3 GB & 0.7 GB & 2.1 GB & 10.5 \\
GLiNER small & 153 M & 0.6 GB & 1.1 GB & 1.7 GB & 12.0 \\
XLM-RoBERTa, fine-tuned & 277 M & 1.1 GB & 0.7 GB & 1.1 GB & 12.0 \\
spaCy, small & --- & 16 MB & 0.3 GB & 0.3 GB & 1.0 \\
TF-IDF + logistic regression & 0.3 M & 4 MB & 0.1 GB & 0.2 GB & 11.0 \\
\midrule
\texttt{llama3.2:3b} & 3.2 B & 2.0 GB & 4.1 GB & --- & --- \\
\texttt{qwen3.5:4b} & 4.2 B & 3.3 GB & 3.7 GB & --- & --- \\
\texttt{phi4-mini} & 3.8 B & 2.5 GB & 4.8 GB & --- & --- \\
\texttt{granite3.3:8b} & 8.2 B & 4.9 GB & 7.9 GB & --- & --- \\
\texttt{gemma4:e4b} & 7.5 B & 6.6 GB & 2.1 GB & --- & --- \\
\texttt{qwen3.5:9b} & 9.0 B & 6.6 GB & 6.7 GB & --- & --- \\
\midrule
Claude Haiku, Sonnet & not published & none & none & --- & --- \\
\bottomrule
\end{tabular}
\caption[What each model needs]{Size and resources on the machine these
comparisons ran on (Intel Core Ultra~7 165U, 14 threads, 64~GB, no GPU
used), which is not the laptop of Section~\ref{sec:setup}: seconds here
do not compare with seconds there. Each model of
the first block was measured in a process of its own over 100 sentences,
above the bare interpreter; cores is processor time over elapsed time.
The second block is Ollama's own report once the model is loaded
(4-bit quantised files); \texttt{gemma4:e4b}'s figure is lower than its
file and is given as reported. Peak memory and cores were not measured
for those models.}
\label{tab:frameworks-resources}
\end{table}

\paragraph{What follows.} For enrichment, a model trained for the task
is the better trade on every axis measured: a classifier trained on the
pipeline's own embeddings for topics, once enough articles are labelled,
and a fine-tuned entity model where products and events are not needed.
A language model, local or hosted, is not needed for either. For
verification the comparison settles nothing at ten claims, and shows
that a hosted model cannot be measured fairly on a public benchmark it
may have read. No phase has a timed human run yet, so the person's
column of this comparison is still empty.
